\section{Metodología de comparación}
\label{sec:metodologia}

La comparación entre ambas implementaciones persigue el tercer objetivo específico, es decir,
medir el tiempo de generación de las pruebas, el uso de recursos de cómputo y de memoria y el costo
de verificación, y caracterizar con esos datos las ventajas y desventajas de cada herramienta.
Esta sección define el enunciado medido, la correspondencia entre las etapas de Noir y de SP1, la
instrumentación empleada, las métricas que se derivan de ella y las condiciones que limitan el
alcance de las mediciones.

\subsection{Enunciado medido y alcance}
\label{sec:metodologia_alcance}

Las mediciones del costo de generación se concentran en la verificación de la firma agregada del comité
de sincronización, es decir, en la comprobación de la Ecuación~\eqref{eq:bls_agregada} sobre la
raíz de firma de un encabezado de la cadena de baliza. Esta verificación constituye el núcleo del
programa de consenso de ambas implementaciones (Secciones~\ref{sec:noir_impl_consenso}
y~\ref{sec:sp1_impl_consenso}) y es, a la vez, su parte más costosa, puesto que exige descomprimir
y sumar claves de la curva BLS12-381, transformar el mensaje en un punto de $\mathbb{G}_2$ y
evaluar emparejamientos. Ambas implementaciones procesaron la misma entrada, correspondiente a una
actualización del comité cuya raíz de firma es
\texttt{0xe4ad1c23\ldots 2a33401a}, y en ambos casos la verificación concluyó con resultado
válido. En Noir, el circuito medido, \texttt{noir\_bls12\_381\_validator}, corresponde a una
versión anterior del programa de consenso, que recibía el punto del mensaje y la clave pública
como entradas no vinculadas al comité ni al encabezado. Ese circuito no incluía, por lo tanto, la
transformación del mensaje en un punto de $\mathbb{G}_2$ ni la vinculación de las claves con el
compromiso del comité, operaciones que sí realiza el programa de SP1 medido. Los circuitos de
consenso que integran el puente, descritos en la Sección~\ref{sec:resultados_gas}, no se midieron
en el servidor.

Los costos del puente se organizan en tres ejes, ya anunciados en la
Sección~\ref{sec:op_tarifas}. El primero es el gas de ejecución de las transacciones de anclaje y
de acuñación. El segundo es su costo monetario en Optimism Sepolia, que resulta del gas consumido,
del precio efectivo del gas y de las tarifas propias de la capa 2. El tercero es el costo de generar
las pruebas fuera de la cadena. Las mediciones de las etapas de generación, descritas en las
subsecciones siguientes, corresponden al tercer eje y se limitan al programa de consenso, ya que
las etapas del programa de inclusión no se midieron con la misma instrumentación. El primer eje se
cubre con la medición del gas de los contratos verificadores de ambas herramientas, de las
transacciones completas de anclaje y de acuñación de Noir y de la lógica de anclaje de SP1
(Sección~\ref{sec:metodologia_gas}). En SP1 no se midió una transacción completa de anclaje ni de
acuñación. El segundo eje no se midió.

\subsection{Correspondencia entre etapas}
\label{sec:metodologia_etapas}

Noir y SP1 recorren etapas equivalentes, aunque las distribuyen de forma distinta entre las
herramientas que las ejecutan (Secciones~\ref{sec:noir_flujo} y~\ref{sec:sp1_flujo}). En ambos
casos, el programa se compila, se ejecuta sobre la entrada para obtener el testigo o la traza, se
deriva la clave de verificación, se genera la prueba y se la verifica. La
Tabla~\ref{tab:metodologia_etapas} recoge la orden ejecutada en cada etapa y la frecuencia con que
debe repetirse durante la operación del puente. Esta frecuencia es la que permite distinguir los
costos fijos, que se pagan una vez por versión del programa, de los costos que se pagan en cada
prueba.

\begin{table}[htbp]
	\centering
	\caption{Correspondencia entre las etapas medidas en Noir y en SP1.}
	\label{tab:metodologia_etapas}
	\small
	\begin{tabular}{>{\raggedright\arraybackslash}p{2.6cm} >{\raggedright\arraybackslash}p{4.3cm} >{\raggedright\arraybackslash}p{4.3cm} >{\raggedright\arraybackslash}p{2.6cm}}
		\toprule
		\textbf{Etapa} & \textbf{Noir} & \textbf{SP1} & \textbf{Frecuencia} \\
		\midrule
		Compilación & \texttt{nargo compile} & \texttt{cargo prove build} & Una vez por versión
		del programa \\
		\addlinespace
		Testigo o ejecución & \texttt{nargo execute} & Anfitrión con \texttt{-{}-execute} & En cada
		prueba \\
		\addlinespace
		Clave de verificación & \texttt{bb write\_vk} & Anfitrión \texttt{vkey} & Una vez por
		versión del programa \\
		\addlinespace
		Prueba & \texttt{bb prove} & Anfitrión con \texttt{-{}-prove} & En cada prueba \\
		\addlinespace
		Verificación & \texttt{bb verify} & Incluida en la etapa de prueba & En cada prueba \\
		\bottomrule
	\end{tabular}
\end{table}

La correspondencia presenta dos asimetrías que deben tenerse en cuenta al interpretar los
resultados. En SP1, la etapa de prueba ejecuta en una sola orden la ejecución del programa, la
generación de la prueba y su verificación local, de modo que su tiempo incluye trabajo que en Noir
se distribuye entre tres etapas. La etapa de ejecución de SP1 se midió por separado porque informa
el número de ciclos del programa, pero no forma parte del costo por prueba, ya que la etapa de
prueba la repite internamente. En Noir, en cambio, el testigo que produce \texttt{nargo execute}
es una entrada necesaria de \texttt{bb prove}, por lo que su tiempo sí integra el costo de cada
prueba.

\subsection{Instrumentación y entorno}
\label{sec:metodologia_instrumentacion}

Cada etapa se ejecutó bajo la orden \texttt{/usr/bin/time -v}, que registra al finalizar el
proceso el tiempo transcurrido de reloj, el tiempo de procesador en modo usuario y en modo
sistema, el porcentaje de uso del procesador, la memoria residente máxima, los fallos de página,
las operaciones de entrada desde el sistema de archivos y el código de salida. El porcentaje de uso
del procesador resulta del cociente entre el tiempo de procesador y el tiempo de reloj, de modo que
un valor superior al 100\,\% indica que el proceso empleó varios hilos en paralelo. La memoria
residente máxima se informa en gibibytes (GiB), equivalentes a $2^{30}$ bytes. Las etapas también
se ejecutaron bajo \texttt{strace}, pero sus registros solo contienen las llamadas al sistema del
intérprete de órdenes que lanzó cada proceso y no las del proceso medido, por lo que se descartaron.

Ambas implementaciones se midieron en el mismo servidor, una máquina virtual de Oracle Cloud
Infrastructure de la forma \texttt{VM.Standard3.Flex} con sistema operativo Linux, configurada con
7 OCPU y 90 GB de memoria. En esta forma, cada OCPU corresponde a un núcleo físico del procesador
con dos hilos de ejecución~\cite{ocicomputeshapes}, de modo que el servidor disponía de 14 hilos y
el porcentaje de uso del procesador podía alcanzar, como máximo, el 1400\,\%. El servidor contaba
además con 150 GB de almacenamiento dedicados al área de intercambio, que el sistema operativo
emplea como extensión de la memoria cuando un proceso supera la memoria física disponible. La
generación de las pruebas se realizó exclusivamente en el procesador. En SP1 esto corresponde a la
variable \texttt{SP1\_PROVER=cpu}, y en Noir al comportamiento predeterminado de Barretenberg. Las
órdenes de Barretenberg se ejecutaron con sus opciones predeterminadas, sin indicar de forma
explícita el objetivo de verificación de la prueba. Las mediciones de SP1 se obtuvieron
entre el 21 y el 22 de julio de 2026, con la cadena de compilación de Succinct basada en
\texttt{rustc} 1.94.0, y las de Noir concluyeron el 10 de agosto de 2026, con Nargo
1.0.0-beta.22 y Barretenberg 5.0.0-nightly.20260522. Cada etapa se ejecutó una sola vez. SP1 se ejecutó también en el modo
simulado (\texttt{SP1\_PROVER=mock}), que omite la generación real de la prueba y está destinado al
desarrollo. En ese modo, la etapa de prueba terminó con el código de salida 101, porque el
anfitrión intentó verificar una prueba de núcleo vacía, de modo que esas mediciones solo se emplean
como referencia del costo de ejecución sin probador.
\subsection{Medición del costo de verificación}
\label{sec:metodologia_gas}

El costo de verificación se midió como el gas de ejecución de las transacciones que verifican las
pruebas en contratos desplegados en redes locales. Esta magnitud no depende del precio del gas
(Sección~\ref{sec:eth_gas}) e incluye el costo intrínseco de la transacción, de 21\,000 unidades,
y el de sus datos, de modo que cada cifra representa una transacción completa. Las mediciones se
realizaron en noviembre de 2026.

El gas de una misma transacción depende de las reglas de la versión de la EVM que aplica la red. La
actualización Osaka de Ethereum incrementó, mediante la propuesta EIP-7883, el costo del contrato
precompilado MODEXP~\cite{eip7883}, que el verificador de UltraHonk emplea, e introdujo con la
propuesta EIP-7825 un límite de $2^{24}$ unidades de gas por transacción~\cite{eip7825}. Optimism
Sepolia aplica ambas propuestas desde su actualización 19~\cite{oposakal2}, de modo que las reglas
posteriores a Osaka son las que representan la cadena de destino. Las mediciones de Noir se
realizaron con esas reglas, en una red de Anvil configurada con la bifurcación BPO1, y se
repitieron con las reglas de la versión anterior, Prague, en la red que Hardhat ejecuta dentro de
su propio proceso~\cite{hardhatdocs}, para permitir la comparación entre ambas.

En Noir, se generaron pruebas reales de los dos circuitos de consenso y del circuito de inclusión
con Barretenberg 5.0.0-nightly.20260522 y el destino \lstinline|evm|, que emplea Keccak en la
transcripción de Fiat-Shamir, y se obtuvo el contrato verificador de cada circuito mediante la orden
\texttt{bb write\_solidity\_verifier -t evm} (Sección~\ref{sec:noir_impl_publicas}). Los contratos
se compilaron con solc 0.8.28, con el optimizador activado y el parámetro \texttt{runs} igual a
uno. Se midieron la transacción de anclaje, que verifica las dos pruebas de consenso y registra las
raíces de transacciones y de recibos del bloque, la transacción de acuñación, que verifica la prueba de inclusión y
acuña el activo, y el gas de cada verificación por separado. Las pruebas se generaron en un equipo
distinto del servidor de la Sección~\ref{sec:metodologia_instrumentacion}, un computador Mac con
18 GB de memoria, por lo que sus tiempos de generación no se comparan con los de la
Sección~\ref{sec:resultados_etapas}.

En SP1, se empleó el contrato verificador Groth16 oficial de la versión 6.1.0, que coincide con la
del kit de desarrollo que resuelven las dependencias del programa de consenso~\cite{sp1contracts}, compilado con solc 0.8.24, el optimizador
activado, el parámetro \texttt{runs} igual a uno y la versión \texttt{shanghai} de la EVM, y
desplegado en Anvil 1.3.5~\cite{anvildocs} con las reglas de la bifurcación Cancun. La función
de verificación de ese contrato solo emplea los precompilados de suma, de multiplicación escalar y
de emparejamiento de la curva BN254 y el de SHA-256~\cite{sp1contracts}, cuyos precios no
modificaron las actualizaciones Prague ni Osaka, por lo que las cifras obtenidas con esas reglas
son válidas también para la cadena de destino. Como no se disponía de una prueba Groth16 del programa
de consenso, la verificación se midió con una prueba Groth16 válida de otro programa, publicada por
un proyecto independiente y compatible con esa versión del verificador~\cite{zekofixture}. Esta
sustitución es admisible porque el costo de verificar una prueba Groth16 no depende del programa
probado (Sección~\ref{sec:snarks_starks}) y solo varía con el tamaño de los valores públicos, que
se incorporan mediante su hash y forman parte de los datos de la transacción. Esa variación se
cuantifica en la Sección~\ref{sec:resultados_gas}. La lógica de \texttt{BeaconLightClient} se
midió por separado, con un verificador simulado que acepta cualquier prueba y con los valores
públicos de 320 bytes del programa de consenso, obtenidos de una ejecución sobre un comité
sintético de un solo miembro.

Ambas mediciones incluyeron controles de validez. En SP1, la prueba válida fue aceptada y la misma
prueba con un byte modificado fue rechazada mediante la reversión de la llamada. En Noir, las
pruebas automatizadas del puente comprobaron, además, el rechazo de las transacciones con una
raíz de transacciones o de recibos alterada, con una prueba de firma adulterada, con un compromiso
del comité distinto del registrado, con un par de raíces que no se anclaron juntas y con entradas
públicas de inclusión modificadas, como un destinatario distinto. Las pruebas del circuito de
inclusión comprobaron, por su parte, el rechazo de un bloqueo revertido, de una transacción de
otra posición, alterada o dirigida a otra función o a otra dirección, de un evento o de un emisor
equivocados y de raíces falsas. En SP1, estos controles se
repitieron sobre el verificador oficial ya desplegado en Optimism Sepolia, y se comprobó además en
esa red el gas que cobran los contratos precompilados de la curva BN254, mediante llamadas que no
modifican el estado. Ninguna medición envió transacciones a una red pública, por lo que no se
obtuvo el costo monetario de las transacciones.

\subsection{Métricas derivadas}
\label{sec:metodologia_metricas}

A partir de las mediciones por etapa se derivan dos métricas que permiten comparar ambas
herramientas en condiciones de uso. La primera es el \textbf{costo en frío}, definido como la suma
de los tiempos de todas las etapas, que representa el tiempo necesario para obtener la primera
prueba desde el código fuente. La segunda resulta de separar el costo fijo $t_{\mathrm{fijo}}$,
formado por la compilación y la obtención de la clave de verificación, del costo por prueba
$t_{\mathrm{prueba}}$, formado por las etapas que se repiten en cada prueba según la
Tabla~\ref{tab:metodologia_etapas}. El tiempo total necesario para generar $n$ pruebas con un
mismo programa es, en consecuencia,
\begin{equation}
	\label{eq:costo_total}
	t_{\mathrm{total}}(n) = t_{\mathrm{fijo}} + n \, t_{\mathrm{prueba}}.
\end{equation}
Si una herramienta presenta un costo fijo mayor y un costo por prueba menor que la otra, existe un
número de pruebas a partir del cual su tiempo total resulta inferior. Ese \textbf{punto de
equilibrio} $n^{*}$ se obtiene igualando la Ecuación~\eqref{eq:costo_total} para ambas
herramientas, lo que conduce a
\begin{equation}
	\label{eq:punto_equilibrio}
	n^{*} = \frac{t^{\mathrm{Noir}}_{\mathrm{fijo}} - t^{\mathrm{SP1}}_{\mathrm{fijo}}}
	{t^{\mathrm{SP1}}_{\mathrm{prueba}} - t^{\mathrm{Noir}}_{\mathrm{prueba}}},
\end{equation}
donde los superíndices identifican la herramienta a la que corresponde cada término.

\subsection{Condiciones de validez}
\label{sec:metodologia_validez}

Seis condiciones acotan la comparabilidad de las mediciones. En primer lugar, la compilación de
SP1 se midió con el ejecutable del programa invitado ya generado, según indica el propio registro
de la orden, que informa la finalización de la compilación en 0,15 segundos. El tiempo obtenido
refleja, por lo tanto, la comprobación de que el ejecutable está actualizado y no la compilación
completa del programa. En segundo lugar, la prueba de SP1 medida es la prueba de núcleo, cuyo
tamaño crece con la extensión de la ejecución (Sección~\ref{sec:sp1_flujo}), y no la prueba
Groth16 que verifica el contrato de destino, por lo que su tiempo no incluye la compresión
recursiva ni la envoltura descritas en la Sección~\ref{sec:sp1_impl_integracion}. En tercer lugar,
el programa de SP1 no emplea los precompilados de la máquina virtual para SHA-256 ni para
BLS12-381 (Sección~\ref{sec:sp1_impl}), de modo que la aritmética de la curva se ejecuta como
instrucciones RISC-V ordinarias. En cuarto lugar, las órdenes de Barretenberg no
indican de forma explícita el objetivo de verificación, y los registros no permiten establecer si
la prueba de Noir medida emplea la función de hash que exige el verificador desplegado en la
cadena de destino. En quinto lugar, el circuito de Noir medido en el servidor corresponde a una
versión anterior del programa de consenso, que no transforma el mensaje en un punto de
$\mathbb{G}_2$ ni vincula las claves con el compromiso del comité
(Sección~\ref{sec:metodologia_alcance}), de modo que demuestra un enunciado menor que el del
programa de SP1 y que el de los circuitos que integran el puente. En sexto lugar, cada etapa se
ejecutó una sola vez, por lo que los resultados no permiten estimar la variabilidad de las
mediciones. Estas condiciones no invalidan
la comparación, pero determinan el sentido en que cada resultado debe interpretarse, aspecto que se
retoma en la Sección~\ref{sec:resultados_discusion}.
